\subsection{Rings and modules of finite length}

If an A-module M is of finite length, we shall denote this length by \(\text{long}_{\mathbf{A}}(\mathbf{M})\) or \(\text{long}(\mathbf{M})\) . Recall that every Artinian ring is Noetherian (Algebra, Chapter VIII, § 6, no. 5, Corollary 3 to Proposition 12) and that every finitely generated module over an Artinian ring is of finite length (loc. cit., Corollary 1 to Proposition 12). Moreover, every Artinian integral domain is a field (Algebra, Chapter VIII, § 6, no. 4, Proposition 9).

PROPOSITION 7. Let M be a finitely generated module over a Noetherian ring A. The following properties are equivalent:

\begin{enumerate}
\def\labelenumi{(\alph{enumi})}
\item
  M is of finite length.
\item
  Every ideal \(\mathfrak{p} \in \operatorname{Ass}(\mathbf{M})\) is a maximal ideal of \(\mathbf{A}\) .
\item
  Every ideal \(\mathfrak{p} \in \operatorname{Supp}(\mathbf{M})\) is a maximal ideal of \(\mathbf{A}\) .
\end{enumerate}

Let \((\mathbf{M}_i)_{0\leqslant i\leqslant n}\) be a composition series of M such that, for \(0\leqslant i\leqslant n - 1\) , \(\mathbf{M}_i / \mathbf{M}_{i + 1}\) is isomorphic to \(\mathrm{A} / \mathfrak{p}_i\) , where \(\mathfrak{p}_i\) is prime (§ 1, no. 4, Theorem 1). If M is of finite length, so is each of the A-modules \(\mathrm{A} / \mathfrak{p}_i\) , which implies that each of the rings \(\mathrm{A} / \mathfrak{p}_i\) is Artinian; but as \(\mathrm{A} / \mathfrak{p}_i\) is an integral domain, it is therefore a field, in other words \(\mathfrak{p}_i\) is maximal; we conclude that (a) implies (b) (§ 1, no. 4, Theorem 2). Condition (b) implies (c) by § 1, no. 3, Proposition 7. Finally, if all the ideals of \(\operatorname {Supp}(\mathbf{M})\) are maximal, so are the \(\mathfrak{p}_i\) (§ 1, no. 4, Theorem 2), hence the \(\mathrm{A} / \mathfrak{p}_i\) are simple A-modules and M is of finite length, which completes the proof.

COROLLARY 1. For every module of finite length M over a Noetherian ring A, \(\operatorname{Ass}(\mathbf{M}) = \operatorname{Supp}(\mathbf{M})\) .

Every element of \(\operatorname{Supp}(\mathbf{M})\) is then minimal in \(\operatorname{Supp}(\mathbf{M})\) and the conclusion follows from § 1, no. 3, Corollary 1 to Proposition 7.

COROLLARY 2. Let M be a finitely generated module over a Noetherian ring A and p a prime ideal of A. For \(M_{p}\) to be a non-zero \(A_{p}\) -module of finite length, it is necessary and sufficient that p be a minimal element of Ass(M).

By § 1, no. 2, Corollary to Proposition 5, \(\operatorname{Ass}_{\mathbf{A}_{\mathfrak{p}}}(\mathbf{M}_{\mathfrak{p}})\) is the set of ideals \(q_{\mathfrak{p}}\) , where \(q\) runs through the set of ideals of \(\operatorname{Ass}(\mathbf{M})\) which are contained in \(p\) . On the other hand, \(p_{\mathfrak{p}}\) is the unique maximal ideal of \(\mathbf{A}_{\mathfrak{p}}\) ; by Proposition 7, for \(\mathbf{M}_{\mathfrak{p}}\) to be an \(\mathbf{A}_{\mathfrak{p}}\) -module of finite length, it is necessary and sufficient that no element of \(\operatorname{Ass}(\mathbf{M})\) be strictly contained in \(p\) . On the other hand, for \(\mathbf{M}_{\mathfrak{p}} \neq 0\) , it is necessary and sufficient by definition that \(\mathfrak{p} \in \operatorname{Supp}(\mathbf{M})\) (Chapter II, § 4, no. 4), that is that \(\mathfrak{p}\) contain an element of \(\operatorname{Ass}(\mathbf{M})\) (§ 1, no. 3, Proposition 7). This proves the corollary.

Remark (1). Let M be a finitely generated module over a Noetherian ring A; let \((\mathbf{M}_{i})_{0\leqslant i\leqslant n}\) be a composition series of M such that, for \(0\leqslant i\leqslant n-1\) , \(M_{i}/M_{i+1}\) is isomorphic to \(A/p_{i}\) , where \(p_{i}\) is a prime ideal of A (§ 1, no. 4, Theorem 1). If p is a minimal element of Ass(M), the length \(\text{long}_{A_{p}}(M_{p})\) is equal to the number of indices i such that \(p_{i}=p\) . For the \((\mathbf{M}_{i})_{p}\) form a composition series of \(M_{p}\) and \((\mathbf{M}_{i})_{p}/(\mathbf{M}_{i+1})_{p}\) is isomorphic to \((\mathbf{A}/p_{i})_{p}\) and hence to \(\{0\}\) if \(p_{i}\neq p\) (since p is minimal in the set of \(p_{i}\) by § 1, no. 4, Theorem 2) and to \((\mathbf{A}/p)_{p}\) which is a field, if \(p_{i}=p\) .

PROPOSITION 8. Let M be a module of finite length over a Noetherian ring A.

\begin{enumerate}
\def\labelenumi{(\roman{enumi})}
\item
  There only exists a single primary decomposition of \(\{0\}\) with respect to \(\mathbf{M}\) indexed by \(\operatorname{Ass}(\mathbf{M})\) (necessarily reduced); let \(\{0\} = \bigcap_{\mathfrak{p} \in \operatorname{Ass}(\mathbf{M})} \mathbf{Q}(\mathfrak{p})\) be this decomposition, where \(\mathbf{Q}(\mathfrak{p})\) is \(\mathfrak{p}\) -primary with respect to \(\mathbf{M}\) .
\item
  There exists an integer \(n_0\) such that, for all \(n \geqslant n_0\) and all \(\mathfrak{p} \in \operatorname{Ass}(\mathbf{M})\) , \(\mathbf{Q}(\mathfrak{p}) = \mathfrak{p}^n\mathbf{M}\) .
\item
  For all \(\mathfrak{p} \in \operatorname{Ass}(\mathbf{M})\) , the canonical mapping of \(\mathbf{M}\) to \(\mathbf{M}_{\mathfrak{p}}\) is surjective and its kernel is \(\mathbf{Q}(\mathfrak{p})\) .
\item
  The canonical injection of \(\mathbf{M}\) into \(\bigoplus_{\mathfrak{p} \in \mathbf{Ass}(\mathbf{M})} (\mathbf{M} / \mathbf{Q}(\mathfrak{p}))\) is bijective.
\end{enumerate}

As every element \(\mathfrak{p} \in \operatorname{Ass}(M)\) is minimal in \(\operatorname{Ass}(M)\) (Proposition 7), assertion (i) follows from no. 3, Proposition 5. As \(M\) is finitely generated, there exists \(n_0\) such that \(\mathfrak{p}^n M \subset Q(\mathfrak{p})\) for all \(\mathfrak{p} \in \operatorname{Ass}(M)\) and all \(n \geqslant n_0\) (no. 1, Remark); but as \(\mathfrak{p}\) is a maximal ideal, \(\mathfrak{p}^n M\) is \(\mathfrak{p}\) -primary with respect to \(M\) (no. 1, Examples 2 and 3) and, as \(\bigcap_{\mathfrak{p} \in \operatorname{Ass}(M)} \mathfrak{p}^n M = \{0\}\) , it follows from (i) that necessarily \(\mathfrak{p}^n M = Q(\mathfrak{p})\) for all \(\mathfrak{p} \in \operatorname{Ass}(M)\) ; whence (ii). As the \(\mathfrak{p}^n\) , for \(\mathfrak{p} \in \operatorname{Ass}(M)\) , are relatively prime in pairs (Chapter II, § 1, no. 2, Proposition 3), the canonical mapping \(M \to \bigoplus_{\mathfrak{p} \in \operatorname{Ass}(M)} (M / \mathfrak{p}^n M)\) is surjective (Chapter II, § 1, no. 2, Proposition 6), whence (iv). Then \(\operatorname{Ass}(Q(\mathfrak{p})) = \operatorname{Ass}(M) - \{\mathfrak{p}\}\) and \(\operatorname{Ass}(M / Q(\mathfrak{p})) = \{\mathfrak{p}\}\) (no. 3, Proposition 4); as the elements of \(\operatorname{Ass}(M)\) are maximal ideals, \(\mathfrak{p}\) is the only element of \(\operatorname{Ass}(M)\) which does not meet \(A - \mathfrak{p}\) ; \(Q(\mathfrak{p})\) is therefore the kernel of the canonical mapping \(j: M \to M_{\mathfrak{p}}\) (§ 1, no. 2, Proposition 6). If \(s \in A - \mathfrak{p}\) , the homothety of \(M / Q(\mathfrak{p})\) with ratio \(s\) is injective by virtue of the relation \(\operatorname{Ass}(M / Q(\mathfrak{p})) = \{\mathfrak{p}\}\) (no. 1, Proposition 1); since \(M / Q(\mathfrak{p})\) is Artinian, this homothety is bijective (Algebra, Chapter VIII, § 1, no. 2, Lemma 3). The canonical mapping \(M \to M / Q(\mathfrak{p})\) is then written \(f \circ j\) , where \(f: M_{\mathfrak{p}} \to M / Q(\mathfrak{p})\) is an A-homomorphism (Chapter II, § 2, no. 2, Proposition 3); as \(\operatorname{Ker}(j) = \operatorname{Ker}(f \circ j) = Q(\mathfrak{p})\) , \(f\) is injective; we conclude that \(j\) is surjective and \(f\) bijective.

COROLLARY. If M is a module of finite length over a Noetherian ring A, then

\[
\operatorname{long} _ {\mathbf {A}} (\mathbf {M}) = \sum_ {\mathfrak {p} \in \mathbf {A s s} (\mathbf {M})} \operatorname{long} _ {\mathbf {A p}} (\mathbf {M} _ {\mathfrak {p}}).\tag{4}
\]

This will follow from Proposition 8 (iv) if we prove that

\[
\operatorname{long} _ {\mathbf {A}} (\mathbf {M} / \mathbf {Q} (\mathfrak {p})) = \operatorname{long} _ {\mathbf {A} _ {\mathfrak {p}}} (\mathbf {M} _ {\mathfrak {p}}).
\]

Now, it follows from Proposition 1 of no. 1 that for all \(s \in \mathbf{A} - \mathfrak{p}\) the homothety with ratio \(s\) on \(\mathrm{M} / \mathrm{Q}(\mathfrak{p})\) is injective; the homothety with ratio \(s\) on every submodule \(\mathbf{R}\) of \(\mathrm{M} / \mathrm{Q}(\mathfrak{p})\) is therefore injective and, as \(\mathbf{R}\) is Artinian, it is bijective (Algebra, Chapter VIII, §2, no. 2, Lemma 3); we conclude that the sub- \(\mathbf{A}\) -modules of \(\mathrm{M} / \mathrm{Q}(\mathfrak{p})\) are the images under the bijection \(f: \mathbf{M}_{\mathfrak{p}} \to \mathbf{M} / \mathbf{Q}(\mathfrak{p})\) of the sub- \(\mathbf{A}_{\mathfrak{p}}\) -modules of \(\mathbf{M}_{\mathfrak{p}}\) (Chapter II, §2, no. 3), whence our assertion.

PROPOSITION 9. Let A be a Noetherian ring. The following conditions are equivalent:

\begin{enumerate}
\def\labelenumi{(\alph{enumi})}
\item
  A is Artinian.
\item
  All the prime ideals of \(\mathbf{A}\) are maximal ideals.
\item
  All the elements of Ass(A) are maximal ideals.
\end{enumerate}

If these conditions are fulfilled, A has only a finite number of prime ideals, which are all maximal and associated with the A-module A; further, A is a semi-local ring and its Jacobson radical is nilpotent.

To say that A is Artinian is equivalent to saying that A is an A-module of finite length; hence (a) and (c) are equivalent by Proposition 7. Clearly (b) implies (c). Finally, (a) implies (b) since every Artinian integral domain is a field. The properties (a), (b) and (c) are therefore equivalent.

Suppose they hold. As every prime ideal of A belongs to Supp(A) and every element of Supp(A) contains an element of Ass(A) (§ 1, no. 3, Proposition 7), it follows from (c) that Ass(A) is the set of all prime ideals of A; then A has only a finite number of prime ideals, all of them maximal and associated with the A-module A. This obviously implies that A is semi-local; finally, we know that the Jacobson radical of an Artinian ring is nilpotent (Algebra, Chapter VIII, § 6, no. 4, Theorem 3).

Remark (2) The conditions of Proposition 9 for a Noetherian ring A imply that the spectrum of A is finite and discrete, every point of \(\operatorname{Spec}(A)\) being therefore closed (Chapter II, §4, no. 3, Corollary 6 to Proposition 11). Conversely, for a Noetherian ring A, to say that every point of \(\operatorname{Spec}(A)\) is closed means that every prime ideal of A is maximal (loc. cit.) and hence this condition is equivalent to that of Proposition 9.

COROLLARY 1. Every Artinian ring A is isomorphic to the direct composition of a finite family of Artinian local rings.

It follows from Proposition 9 and Proposition 8 (iii) and (iv) that, if \((\mathfrak{m}_i)_{1\leqslant i\leqslant n}\) is the family of maximal ideals of A, the canonical mapping \(\mathbf{A}\to \prod_{i}\mathbf{A}_{m_i}\) is bijective.

Remark (3). This corollary can also be deduced from the fact that \(\operatorname{Spec}(A)\) is finite and discrete and Chapter II, § 4, no. 3, Proposition 15.

COROLLARY 2. Let A be a Noetherian ring and m an ideal of A. The following conditions are equivalent:

\begin{enumerate}
\def\labelenumi{(\alph{enumi})}
\item
  A is a semi-local ring and m is a defining ideal of A.
\item
  A is a Zariski ring with the m-adic topology and A/m is Artinian.
\end{enumerate}

If (a) holds, A is a Zariski ring with the m-adic topology (Chapter III, § 3, no. 3, Example 3); further, as by hypothesis m contains a power of the Jacobson radical r of A, every prime ideal of A which contains m also contains r (Chapter II, § 1, no. 1, Proposition 1); it is therefore maximal, since r is a finite intersection of maximal ideals (loc. cit.); Proposition 9 then shows that A/m is Artinian. Conversely, if (b) holds, every maximal ideal p of A contains the Jacobson radical of A and hence contains m (Chapter III, § 3, no. 3, Proposition 6); as A/m is Artinian, the ideals p/m are finite in number (Proposition 9) and hence A has only a finite number of maximal ideals, which implies that it is semi-local.

COROLLARY 3. Let A, A' be two rings and h a homomorphism from A to A'. Suppose that A is semi-local and Noetherian and that A' is a finitely generated A-module. Then the ring A' is semi-local and Noetherian; if m is a defining ideal of A, mA' is a defining ideal of A'.

We know that A' is a Zariski ring with the mA'-adic topology (Chapter III, § 3, no. 3, Proposition 7). As A/m is Artinian (Corollary 2) and A'/mA' is a finitely generated (A/m)-module, A'/mA' is an Artinian ring, hence A' is semi-local and mA' is a defining ideal of A' (Corollary 2).

COROLLARY 4. Let A be a complete semi-local Noetherian ring, m a defining ideal of A, E a finitely generated A-module and \((\mathbf{F}_{n})\) a decreasing sequence of submodules of E such that \(\bigcap_{n} F_{n} = 0\) . Then, for all p \textgreater{} 0, there exists n \textgreater{} 0 such that \(F_{n} \subset m^{p}E\) .

As A is a Zariski ring, E is Hausdorff and the \(F_{n}\) are closed under the m-adic topology. On the other hand, E is complete (Chapter III, § 2, no. 12, Corollary 1 to Proposition 16). Finally, \(E/m^{p}E\) is a finitely generated module over the ring \(A/m^{p}\) , which is Artinian (Corollary 2); then \(E/m^{p}E\) is an Artinian \((A/m^{p})\) -module and hence an Artinian A-module. The corollary then follows from Chapter III, § 2, no. 7, Proposition 8.

COROLLARY 5. In a complete semi-local Noetherian ring every decreasing sequence of ideals whose intersection is 0 is a filter base which converges to 0.

It is sufficient to apply Corollary 4 to the A-module A.

PROPOSITION 10. Let A be a Noetherian ring and \(\mathfrak{p}_1, \ldots, \mathfrak{p}_n\) the prime ideals associated with the A-module A, where \(\mathfrak{p}_i \neq \mathfrak{p}_j\) for \(i \neq j\) .

\begin{enumerate}
\def\labelenumi{(\roman{enumi})}
\item
  The set \(\mathbf{S} = \bigcap_{i=1}^{n} (\mathbf{A} - \mathfrak{p}_i)\) is the set of elements which are not divisors of 0 in A.\\
\item
  If all the \(\mathfrak{p}_i\) are minimal elements of \(\operatorname{Ass}(\mathbf{A})\) , the total ring of fractions \(\mathbf{S}^{-1}\mathbf{A}\) of \(\mathbf{A}\) is Artinian.
\item
  If the ring A is reduced, all the \(p_{i}\) are minimal elements of Ass(A) (and therefore are the minimal elements of Spec(A)) and each of the \(A_{p_{i}}\) is a field; for each index i, the canonical homomorphism \(S^{-1}A \to A_{p_{i}}\) (Chapter II, §2, no. 1, Corollary 1 to Proposition 2) is surjective and its kernel is \(S^{-1}p_{i}\) ; finally the canonical homomorphism from \(S^{-1}A\) to \(\prod_{i=1}^{n} (S^{-1}A/S^{-1}p_{i})\) is bijective.
\end{enumerate}

The fact that S is the set of elements which are not divisors of 0 in A has already been seen ( \(\S\) 1, no. 1, Corollary 3 to Proposition 2). The prime ideals of \(S^{-1}A\) are

of the form \(\mathbf{S}^{-1}\mathfrak{p}\) , where \(\mathfrak{p}\) is a prime ideal of \(\mathbf{A}\) contained in \(\bigcup_{i=1}^{n} \mathfrak{p}_i\) (Chapter II, §2, no. 5, Proposition 10) that is contained in one of the \(\mathfrak{p}_i\) (Chapter II, §1, no. 1, Proposition 2). If \(\mathfrak{p}_i\) is a minimal element of \(\operatorname{Ass}(\mathbf{A})\) , it is a minimal element of \(\operatorname{Spec}(\mathbf{A})\) (§1, no. 3, Corollary to Proposition 7); if each of the \(\mathfrak{p}_i\) is a minimal element of \(\operatorname{Ass}(\mathbf{A})\) , we then see that the prime ideals of \(\mathbf{S}^{-1}\mathbf{A}\) are the \(\mathbf{S}^{-1}\mathfrak{p}_i\) and they are therefore all maximal, which proves that \(\mathbf{S}^{-1}\mathbf{A}\) is Artinian (Proposition 9).

Suppose finally that the ring A is reduced. Then \(\bigcap_{i=1}^{n} p_i = \{0\}\) (§ 1, no. 3,

Corollary 2 to Proposition 7). We deduce that \(\{0\} = \bigcap_{i=1}^{n} p_i\) is a reduced primary decomposition of the ideal \(\{0\}\) (no. 3, Corollary to Proposition 4); in particular, none of the \(p_i\) can contain a \(p_j\) of index \(j \neq i\) and therefore the \(p_i\) are all minimal elements of \(\operatorname{Ass}(A)\) . The ring \(S^{-1}A\) is then Artinian by (ii). The \(S^{-1}p_i\) are prime ideals associated with the \(S^{-1}A\) -module \(S^{-1}A\) ( \(\S\) 1, no. 2, Corollary to Proposition 5) and \(\{0\} = S^{-1}\left(\bigcap_{i=1}^{n} p_i\right) = \bigcap_{i=1}^{n} S^{-1}p_i\) (Chapter II, \(\S\) 2, no. 4); as the \(S^{-1}p_i\) are distinct, \((S^{-1}p_i)_{1 \leqslant i \leqslant n}\) is a reduced primary decomposition of \(\{0\}\) in \(S^{-1}A\) (no. 3, Corollary to Proposition 4). Proposition 8 then shows that the canonical homomorphism \(g_{i} \colon S^{-1}\mathbf{A} \to (\mathbf{S}^{-1}\mathbf{A})_{p_{i}}\) is surjective and has kernel \(S^{-1}p_{i}\) and the canonical homomorphism \(S^{-1}\mathbf{A} \to \prod_{i=1}^{n} (\mathbf{S}^{-1}\mathbf{A}/\mathbf{S}^{-1}p_{i})\) is bijective. We know moreover that the canonical homomorphism \(S^{-1}\mathbf{A} \to A_{p_{i}}\) is composed of \(g_{i}\) and an isomorphism \((\mathbf{S}^{-1}\mathbf{A})_{\mathbf{S}^{-1}p_{i}} \to A_{p_{i}}\) (Chapter II, § 2, no. 3, Proposition 7). Finally, it follows from Proposition 8 that \((\mathbf{S}^{-1}\mathbf{A})_{\mathbf{S}^{-1}p_{i}}\) is isomorphic to \(S^{-1}\mathbf{A}/\mathbf{S}^{-1}p_{i}\) and hence is a field since \(S^{-1}p_{i}\) is a maximal ideal.

